\section{Relative programs for rational nef line data}\label{sec:relative}

Rational line bundles on projective modifications provide the relative
constructions used in the contraction induction and the companion
flip-termination theorem. We establish ordinary representatives over Stein
neighborhoods, elementary relative steps, a finite auxiliary program,
and extraction of prescribed divisors. All these constructions take
place over a fixed compact base.  Their numerical spaces use degrees
of global line bundles, rather than the full Bott--Chern space of the
ordinary contractions above.

The nef data in this section are rational line bundles, distinct from the
possibly transcendental nef classes used in the supporting-contraction proof.

\subsection{Rational bundles and generalized pairs}

\begin{convention}\label{conv:rational-bundles}
A rational line bundle is an element of $\Pic(T)\otimes_{\Z}\Q$.
An equality of rational line bundles means an isomorphism of holomorphic
line bundles after a positive common multiple.  A rank-one reflexive
sheaf is rationally invertible if one of its positive reflexive powers
is invertible.  We use additive notation for these objects, their
tensor products, and their reflexive transforms.  In particular, $K_T$
denotes the reflexive canonical sheaf, without a choice of a global
meromorphic canonical form.

For a projective morphism $T\to V$, the words relatively nef and
relatively ample, applied to rational line bundles, refer respectively
to degrees on curves in the fibers and ampleness of a positive integral
power.  Real combinations of rational line bundles have the corresponding
numerical meanings.  Absolute nefness of a bundle on a compact K\"ahler
space implies relative nefness in this sense.
\end{convention}

\begin{definition}\label{def:generalized-data}
Rational generalized b-line data on a normal compact space $T$ consist
of an effective rational divisor $B$, a projective modification
$p\colon W\to T$, and a rational line bundle $M_W$ on $W$.
The bundle is pulled back on every higher model; this b-object is denoted
by $\bM$.  We require $M_W$ to be nef, or to be nef over a specified
base in a relative argument.  Its trace on $T$ is obtained by taking a
proper direct image of a cleared power and then its double dual.  The
adjoint
\[
 D_T=K_T+B+M_T
\]
is required to be rationally invertible.  We may increase $W$ to a smooth
projective log resolution carrying the data.  On such a model the
boundary $B_W$ is determined by the crepant formula
\begin{equation}\label{eq:generalized-crepant}
 K_W+B_W+M_W=p^*D_T.
\end{equation}
The equality uses the natural meromorphic identifications over the
isomorphism locus.  Log discrepancy is
$a(E;T,B+\bM)=1-\operatorname{coeff}_E B_W$, evaluated on a model
carrying $E$.  Generalized klt and generalized lc, abbreviated gklt and
glc, mean that all these discrepancies are positive and nonnegative,
respectively.  A generalized terminal pair is gklt and has
$a(E;T,B+\bM)>1$ for every exceptional prime divisor over $T$.

A generalized dlt pair, abbreviated gdlt, is a glc pair for which there
is a Zariski open set $T^\circ\subset T$ such that
$(T^\circ,B|_{T^\circ})$ is simple normal crossing, the b-data descend
there, and $T^\circ$ contains the general point of every
zero-discrepancy center.  We choose the carrier and subsequent log
resolutions to be isomorphisms on this open set.  Thus the centers in
$T^\circ$ are strata of the coefficient-one boundary.
\end{definition}

We also allow real combinations of a fixed finite set of rational data
for local cone arguments and real scaling. The extraction and termination
statements below use rational
data.

\begin{lemma}[Traces and exceptional comparisons]\label{lem:trace}
Let $p\colon W\to T$ be a projective modification of normal compact
spaces.
\begin{enumerate}[label=\textup{(\roman*)}]
\item For a line bundle $L$ on $W$, $(p_*L)^{**}$ is coherent and
rank one.  If its appropriate reflexive power is invertible, the natural
comparison between that power pulled back to $W$ and the corresponding
power of $L$ is an exceptional Cartier divisor.  This assertion uses
actual bundles and is compatible with further pullback.
\item Rank-one reflexive sheaves can be transported coherently across
a bimeromorphic map by pullback, proper direct image, and double dual on
a common projective resolution.  Their transforms agree on the common
codimension-one open set and are independent of the chosen resolution
when the map is small.
If $L=L'+G$ is an identity of rational bundles on the carrier with
$G$ a rational divisor, their reflexive traces satisfy
$L_T=L'_T+p_*G$.  The trace of $K_W$, or of its pullback to a higher
model, is $K_T$.
\item If $T$ is globally strongly $\Q$-factorial and $\bM$ is nef
over a base over which $p$ is projective, then
\begin{equation}\label{eq:nef-trace-defect}
 F_M:=p^*M_T-M_W\geq0
\end{equation}
is exceptional.  Consequently, forgetting the nef part of a gklt,
respectively gdlt, pair on $T$ gives an ordinary klt, respectively dlt,
pair.  A gdlt pair becomes gklt after replacing $B$ by
$B-\epsilon\lfloor B\rfloor$, for every sufficiently small
$\epsilon>0$.
\end{enumerate}
\end{lemma}

\begin{proof}
Proper direct image preserves coherence.  On the isomorphism locus its
rank is one, and double dual gives a rank-one reflexive sheaf on the
normal target.  Locally on $T$, choose meromorphic generators of this
coherent sheaf and of a cleared invertible power.  Pulling back their
comparison gives a meromorphic comparison with the bundle upstairs.
One can see that the comparison is meromorphic, rather than merely
defined on a punctured open set, by using a local finite presentation
of $p_*L$: its generators evaluate to holomorphic sections of $L$, and
their ratios supply local denominators.  The comparisons agree on
overlaps because they agree on the isomorphism locus.  Their divisor
is exceptional.  This also proves compatibility with powers and
pullback.  A global meromorphic frame was not used.

For a common resolution $W$ with maps to the two spaces, first pull
back the sheaf, remove torsion, and take its double dual; after an
additional projective principalization it is a line bundle.  Proper
direct image to the other space followed by double dual is coherent.
On a common codimension-one open set it is the required transform.
Normality gives uniqueness of a reflexive extension across codimension
at least two, proving the assertion for small maps.
The additive trace identity is checked at every codimension-one point
of $T$, where the modification is an isomorphism, and then extended
reflexively.  The same observation applies to canonical sheaves and
their pullbacks.  It requires compatible local canonical forms only.

Global strong $\Q$-factoriality makes $M_T$ rationally invertible, so
the first assertion defines $F_M$.  Its negative is $p$-nef and has
zero pushforward.  Lemma~\ref{lem:negativity} gives $F_M\geq0$.  The ordinary crepant boundary is $B_W-F_M$:
\eqref{eq:generalized-crepant} gives explicitly
\[
 p^*(K_T+B)-K_W=B_W+M_W-p^*M_T=B_W-F_M.
\]
Thus ordinary log discrepancies are at least the generalized ones.
Every ordinary zero center is a generalized zero center and is
therefore generically in the specified simple normal crossing open
set.  This proves the klt and dlt assertions.

The effective divisor $\lfloor B\rfloor$ is rational Cartier on $T$.
Decreasing it subtracts the effective divisor
$\epsilon p^*\lfloor B\rfloor$ from the crepant boundary.  A former
positive discrepancy stays positive.  Every former zero-discrepancy
place has center in the floor, since its center is generically a
coefficient-one stratum, and hence has strictly positive order on the
pullback of the floor.  Its new discrepancy is positive as well.  The
boundary on $T$ remains effective for $0<\epsilon<1$.
\end{proof}

\subsection{Ordinary representatives over Stein neighborhoods}

We use the following local analytic inputs with their stated domains.
A Stein compact has property \textup{(P)} if its intersections with
analytic sets defined near the compact have finitely many connected
components.
For a projective morphism over a Stein compact satisfying property
\textup{(P)}, the ordinary klt cone theorem has finitely many rays in
each region truncated by a relatively ample divisor
\cite[Theorem~7.2]{Fujino2022}.  The base point free theorem generates
every sufficiently large integral power of a nef line bundle $L$ when
$aL-(K+\Delta)$ is relatively ample
\cite[Theorems~6.2 and~6.5]{Fujino2022}.  Ordinary klt flips exist for
small projective morphisms of normal analytic spaces without a
$\Q$-factoriality assumption \cite[Theorem~1.14]{Fujino2022}.  Finally, on
a smooth source the multigraded adjoint ring with simultaneous SNC
subunit rational boundaries and a common relatively ample rational
part is locally finitely generated
\cite[Theorem~3.1]{DHP2024}.  We explain the passages from these local
statements to our data and compact base.

\begin{lemma}[Local ordinary representatives]\label{lem:ordinary-representative}
Let $\pi\colon T\to V$ be projective and bimeromorphic, let
$(T,B+\bM)$ be gklt rational data, and suppose the carrier is
projective over $T$ and $M_W$ is nef over $V$.  Let $H$ be a
$\pi$-ample rational line bundle.  Near any sufficiently small Stein
compact in $V$ there are an effective ordinary klt boundary $\Delta$,
an effective rational divisor $H_0\simq H$, and $\eta>0$ such that
\begin{equation}\label{eq:ordinary-replacement}
 K_T+\Delta\simq D_T,\qquad \Delta\geq\eta H_0.
\end{equation}
These are actual rational line-bundle equivalences on that
neighborhood.  For $0\leq\delta\leq\eta/2$,
$\Delta-\delta H_0$ is effective klt and represents $D_T-\delta H$.
For finitely many rational adjoints of this type on a common initial
space, the representatives can have a common positive ample part.
Their convex combinations have the same properties.  The local cone
conclusions also hold for gklt real combinations of finitely many
rational data.
\end{lemma}

\begin{proof}
Write $p\colon W\to T$ for a smooth carrier, with $p$ projective, and put
$\rho=\pi p$.  Work on a Stein neighborhood $U$ of the prescribed
compact.  Choose a $\rho$-ample rational line bundle $A$.  For an
integer $m$ clearing denominators, the sheaf
\[
 \rho_*\cO_W\bigl(m(M_W-A)\bigr)
\]
is coherent and has rank one over the bimeromorphic isomorphism locus.
It is nonzero.  Cartan's Theorem~A gives a nonzero section over $U$,
whose pullback defines an effective divisor.  Hence, after division
by $m$,
\begin{equation}\label{eq:relative-kodaira-section}
 M_W\simq A+E,\qquad E\geq0.
\end{equation}
This proves relative bigness directly, including when $M_W=0$.
The divisor $E$ need not be exceptional.  The equivalence comes from
a section of the specified bundle, and introduces no unspecified
numerical or base twist.

Keep $W$ and the $\rho$-ample bundle $A$ fixed.  Choose an auxiliary
log resolution of the finite fixed supports only to test the
singularities.  On that auxiliary resolution there are finitely many
crepant coefficients over the compact preimage, all strictly below
one.  Hence a sufficiently small rational $\epsilon>0$ makes
$(W,B_W+\epsilon E)$ sub-klt.
The rational bundle
\[
 A_\epsilon=(1-\epsilon)M_W+\epsilon A
\]
is $\rho$-ample.  Choose a rational $\eta>0$ such that
$A_\epsilon-\eta p^*H$ is still $\rho$-ample.  On a smaller
neighborhood, sufficiently large powers of these relatively ample
bundles are relatively generated.  Cartan's Theorem~A supplies finitely
many sections generating over the compact preimage; shrink once more
so that they generate throughout.  Divided general members give
\[
 H_0\geq0,\quad H_0\simq H,
 \qquad G\geq0,\quad
 G\simq A_\epsilon-\eta p^*H.
\]
Choose the members on $T$ and $W$, respectively, with large
denominators, and test generality on the auxiliary log resolution.
Their pulled-back systems remain free; we do not claim that a pulled-back
ample bundle stays ample.  Analytic
Bertini gives
\[
 C_W=B_W+\epsilon E+G+\eta p^*H_0
\]
sub-klt.  Here generation is used for the ample systems only; the
nonzero section producing $E$ was not claimed to generate its bundle.
The analytic generation and Bertini statements used in this step are
\cite[Theorems~2.12 and~2.20]{DHP2024}.

Push down to obtain
\[
 \Delta=B+\epsilon p_*E+p_*G+\eta H_0\geq\eta H_0.
\]
Only exceptional coefficients of $B_W$ could be negative, so
$\Delta$ is effective.  The actual upstairs identity is
$K_W+C_W\simq p^*D_T$.  Choose an integer clearing this identity,
push forward the resulting line-bundle isomorphism, and take double
duals.  Agreement in codimension one gives
\begin{equation}\label{eq:replacement-cartier}
 \cO_T\bigl(m(K_T+\Delta)\bigr)\simeq\cO_T(mD_T).
\end{equation}
In particular, the new ordinary adjoint is rational Cartier.  We have
not inferred Cartierness of newly chosen local divisors from global
strong $\Q$-factoriality.  The divisor
\[
 K_W+C_W-p^*(K_T+\Delta)
\]
is exceptional and $p$-numerically trivial by the same line-bundle
identity.  Lemma~\ref{lem:negativity} in both signs makes it zero.
Thus $C_W$ is the ordinary crepant boundary, and $(T,\Delta)$ is klt.
Subtracting $\delta H_0$ gives the asserted representation of
$D_T-\delta H$ and leaves the common ample part $(\eta/2)H_0$.

For finitely many vertices, use one carrier and one $H_0$, choose their
effective divisors in \eqref{eq:relative-kodaira-section}, and resolve
all fixed supports together.  Take a common sufficiently small
$\epsilon$ and then a common $\eta$.  Choose the finitely many free
members generally.  On the resulting common resolution all the
crepant coefficients are uniformly below one.  Convex combinations
are therefore klt and keep the same ample part.  The parameter space
can be the simplex of vertex weights; no affine independence of the
original vertices is required.

For real data, first express the bundle in a finite real span of
rational bundles on the carrier.  Every signed rational generator has
an effective representative over the Stein neighborhood by the same
Cartan argument.  Thus a real nef bundle also has an expression
$M_W\sim_{\R}A+E$ with $A$ relatively ample and $E\geq0$ in a fixed
finite divisor span.  The preceding small-$\epsilon$ construction
applies.  A real relatively ample bundle can be written as a positive
combination of finitely many nearby rational relatively ample bundles;
divided general members of those bundles give the required real
boundary.  The ordinary real klt cone theorem is consequently
applicable.  Only the rational construction is used for line-bundle
descent and for the finite adjoint rings.
\end{proof}

Lemma~\ref{lem:ordinary-representative} also supplies the hypothesis
needed to apply integral descent to a relatively antiample generalized
klt adjoint.  Indeed its ordinary representative is klt and has the
same relative line-bundle class.  If an integral line bundle $L$ has
zero degree on every contracted curve, then
$aL-(K_T+\Delta)$ is relatively ample.  Lemma~\ref{lem:descent}
therefore gives the actual identity $L=f^*f_*L$; the consecutive-power
argument in its proof introduces no extra integral multiple.  Clearing
one denominator gives the corresponding assertion for rational line
bundles.  The representatives may differ on the Stein neighborhoods:
the descended bundle and its evaluation map are intrinsic and glue.

\subsection{Global relative degrees and elementary steps}

For a projective morphism $\pi\colon T\to V$ of compact spaces, let
$N^1_{\mathrm{gl}}(T/V)_{\Q}$ be the space of global rational line
bundles modulo equality of degrees on all curves contracted by $\pi$.
Put $N^1_{\mathrm{gl}}(T/V)=N^1_{\mathrm{gl}}(T/V)_{\Q}\otimes\R$,
and let $N_{1,\mathrm{gl}}(T/V)$ be the dual space generated by relative
curve functionals.  These spaces are finite-dimensional: pull back to
a compact resolution and use the first Chern class in its finite
dimensional cohomology.  Write
\[
 \overline{\mathrm{NE}}_{\mathrm{gl}}(T/V)
 \subset N_{1,\mathrm{gl}}(T/V)
\]
for the closed cone generated by these functionals.  A relatively ample
global line bundle gives a compact normalized slice of this cone.

Here are details about the compactness and the local cone statements
which will be used.  Choose a finite cover of $V$ by interiors of
Stein compacts $W_\nu$ in Stein open sets $U_\nu$.  We choose them
semianalytic in coordinate charts, so they satisfy property
\textup{(P)}.  On each chart the ordinary
relative numerical cone has a compact slice normalized by the
restriction of a fixed global relatively ample bundle $H$.  Restricting
global bundles defines a linear projection from the local curve space
to $N_{1,\mathrm{gl}}(T/V)$.  The image of that slice is compact and
still lies in $H=1$.  The convex hull of these finitely many compact
images is compact and is precisely the global normalized slice.
Indeed every relative curve lies over one of the compacts, and passage
to closure commutes with this finite compact convex hull.

Apply Lemma~\ref{lem:ordinary-representative} on the charts to a gklt
adjoint $J$.  The local cone theorem and projection show that, for each
$c>0$, only finitely many generating curve rays are needed in the
global region $(J+cH)<0$.  The remaining summand is in
$(J+cH)\geq0$.  If $J$ is nef, apply the same argument to the
ordinary representative of $J-\delta H$, with a further ample
perturbation $(\delta/2)H$.  In the slice $H=1$ this gives a finite
set of curve rays near $J=0$ and a remainder on which
$J\geq\delta/2$.  In particular the face $J=0$ is generated by
finitely many curves.  It also gives the following useful openness
statement: if a linear perturbation is positive on that zero face,
it is positive on the whole slice for all sufficiently small positive
parameters.  This follows by checking the finite rays and using the
uniform $\delta/2$ margin on the remainder.

For a rational gklt adjoint, every negative global extremal ray has a
rational supporting line bundle.  To see this, use the finite
decomposition in a truncated negative neighborhood of that ray.
The normalized point of an extremal ray is outside the compact convex
hull of the other generators and the remaining nonnegative part.
Separation gives a nef functional vanishing only on that ray, with a
strict margin on the other part.  These functionals contain an open
subset of its perpendicular hyperplane.  The ray is generated by a
curve, so the hyperplane is rational; choose a rational functional
there.  It is a rational line-bundle class by the definition of
$N^1_{\mathrm{gl}}$.  Subtracting a sufficiently small positive
multiple of the adjoint is strictly positive on the normalized slice.
After rescaling the support, its difference from the adjoint is
relatively ample.  Strict positivity here implies ampleness on each
projective fiber by Kleiman's criterion, and relative ampleness follows
over a neighborhood of each fiber.  This reasoning involves line
bundle degrees on projective fibers only.

\begin{lemma}[Sheaves and signs across a step]\label{lem:sheaf-step}
Let $T$ be globally strongly $\Q$-factorial.  Suppose a projective
bimeromorphic contraction $f\colon T\to Z$ has all its contracted
curves in one nonzero ray for global rational line-bundle degrees.
Let $J$ be a rational line bundle with $-J$ $f$-ample, admitting the
local klt representatives required in Lemma~\ref{lem:descent}.
\begin{enumerate}[label=\textup{(\roman*)}]
\item If $f$ is divisorial, its exceptional locus is a single prime
$E$, with $E\cdot R<0$, and $Z$ is globally strongly $\Q$-factorial.
Every adjoint pushed down to $Z$ is rational Cartier.
\item If $f$ is small and its $J$-positive small projective model
$f^+\colon T^+\to Z$ exists, then $T^+$ is globally strongly
$\Q$-factorial.  Every global rational line bundle $L$ on $T$ has a
rationally invertible transform $L^+$, and
\begin{equation}\label{eq:step-line-adjustment}
 L=f^*A+cJ\quad\Longrightarrow\quad
 L^+=(f^+)^*A+cJ^+
\end{equation}
for an actual rational line bundle $A$ on $Z$, where
$c=(L\cdot R)/(J\cdot R)\in\Q$.  In particular, a negative, zero,
or positive degree for $L$ becomes respectively relative ampleness,
relative triviality, or relative antiampleness for $L^+$.
\end{enumerate}
\end{lemma}

\begin{proof}
For an exceptional prime $E$, global strong $\Q$-factoriality makes $E$
rational Cartier.  If it were $f$-nef,
Lemma~\ref{lem:negativity}\textup{(i)} applied to $-E$ would give
$-E\geq0$.  Thus $E\cdot R<0$.  Every contracted curve
must lie in $E$, since an effective rational Cartier divisor has
nonnegative intersection with a curve not contained in its support.
Every point of a nontrivial projective fiber lies on a curve in that
fiber.  Hence the entire exceptional locus is $E$, and no second
exceptional prime exists.

For a global rank-one reflexive sheaf $\cF$ on $Z$, form its reflexive
pullback to $T$.  A positive power is a line bundle.  Subtract a
rational multiple of $E$ so that its degree is zero on $R$ and apply
Lemma~\ref{lem:descent}.  The descended rational bundle agrees with
$\cF$ to that power on the complement of a codimension-two set in
$Z$.  Reflexivity extends the identity.  This proves the global strong
condition and, in particular, applies to the pushed-forward canonical,
boundary, and nef-trace sheaves.

In the small case, $L-cJ$ has zero degree and descends by
Lemma~\ref{lem:descent}.  The right-hand side of
\eqref{eq:step-line-adjustment} is a rational line bundle on $T^+$
and agrees with the transform of $L$ on the common open set.  It is
therefore the reflexive transform.  The sign conclusions follow from
relative ampleness of $J^+$ and the sign of $c$.

Conversely, transport an arbitrary global rank-one reflexive sheaf
$\cG$ on $T^+$ to $T$ as in Lemma~\ref{lem:trace}.  A reflexive power
of its transform is invertible on $T$.  The preceding argument makes
its transform a rational line bundle on $T^+$.  It agrees with the
corresponding reflexive power of $\cG$ on the common codimension-one
open set, and hence everywhere.  This is global strong $\Q$-factoriality
on $T^+$; no assertion about arbitrary open subsets is involved.
\end{proof}

\begin{definition}\label{def:relative-elementary}
An elementary relative step is a
projective divisorial contraction
or a small projective negative-to-positive adjoint surgery over a
normal compact K\"ahler base, with globally strongly $\Q$-factorial
source and output.  All contracted curves on the negative side span
one nonzero ray when paired with global rational line bundles.  The
boundary is pushed forward or strictly transformed, and the same
b-data are used on common higher models.  The small positive model is
the relative Proj of the actual adjoint algebra.  No condition on its
relative Bott--Chern dimension is included.
\end{definition}

\begin{proposition}[Relative continuation]\label{prop:relative-steps}
Let $\pi\colon T\to V$ be a projective bimeromorphic morphism of
normal compact K\"ahler spaces.  Assume $T$ is globally strongly
$\Q$-factorial and $(T,B+\bM)$ is rational gdlt, with projective
carrier and nef data over $V$.  If its adjoint is not curve-nef over
$V$, there is an elementary negative step over $V$.  Its target is
projective over $V$, compact K\"ahler, globally strongly
$\Q$-factorial, and gdlt with the transformed data.  The construction
is available after every finite prefix whose final adjoint is not
curve-nef over $V$.  For gklt data one may choose
any negative extremal ray in the global degree cone.
\end{proposition}

\begin{proof}
First assume gklt.  The global cone construction above gives a negative
ray and a nef rational supporting line bundle $L$ which vanishes
precisely on that ray and for which a multiple minus the adjoint is
relatively ample.  Apply the ordinary base point free theorem on the
finite Stein cover, using Lemma~\ref{lem:ordinary-representative}.
It generates a common positive power of $L$ relative to $V$.
The relative section morphism and its Stein factorization give
$f\colon T\to Z$ with connected fibers and normal target, contracting
exactly the curves on the selected ray.  The target is projective over
$V$: it is finite over the image in the projectivization of the
coherent sheaf of relative sections.  Since $T\to V$ is bimeromorphic,
$Z$ has the same dimension and $f$ is bimeromorphic.  The negative
adjoint is relatively ample, because on every $f$-fiber its class is
a positive multiple of a relatively ample class in the one-dimensional
global degree space.

If $f$ is divisorial, apply Lemma~\ref{lem:sheaf-step}.  If it is
small, apply Lemma~\ref{lem:ordinary-representative} locally over
$Z$, followed by the ordinary analytic klt flip theorem
\cite[Theorem~1.14]{Fujino2022}.  Equivalently, the actual relative
adjoint algebra is locally finitely generated by
\cite[Corollary~3.7]{DHP2024}; the flip theorem also supplies its
normality and smallness.  The local models glue by the intrinsic
algebra.  Compactness permits a common Veronese generated in degree
one.  On the common codimension-one open set its tautological bundle
is the stated adjoint power.  Both sheaves are reflexive on the normal
positive side, so this identity extends globally.  Thus the positive
adjoint is rational Cartier and relatively ample as an actual bundle.
Lemma~\ref{lem:sheaf-step} gives global strong $\Q$-factoriality.

All spaces constructed are projective over a compact K\"ahler space.
A relatively ample line bundle has a metric with positive curvature
in the fiber directions; adding a sufficiently large pullback of a
K\"ahler form on the base gives a K\"ahler form.  Thus the spaces are
compact K\"ahler.  The common carrier remains projective over the
models: take the main component of the graph over $Z$ and resolve it.
The pulled-back nef data are still nef over $V$.

For gdlt data, replace $B$ by
$B-\epsilon\lfloor B\rfloor$ with small rational $\epsilon>0$.
By Lemma~\ref{lem:trace} this is gklt; denote its adjoint by $D_\epsilon$.
Normalize by a relatively ample bundle $H$, and choose a relative
curve class $z$ with $H\cdot z=1$ and $D\cdot z=-a<0$.
Choose $\epsilon$ so that
$|(D_\epsilon-D)\cdot u|<a/4$ on the whole compact normalized slice.
Choose an extreme point $u_0$ of the face where $D$ attains its
minimum on this slice.  It is an extreme point of the entire slice,
so it spans a global extremal ray $R$.  Since
$D\cdot u_0\leq D\cdot z=-a$, the uniform perturbation bound gives
$D_\epsilon\cdot u_0<-3a/4$.  Thus $R$ is negative for both
adjoints.  The gklt construction contracts this global extremal ray
for the perturbed pair.
The original adjoint has the same negative sign on the ray.  In the
small case Lemma~\ref{lem:sheaf-step} identifies its positive transform
with a positive multiple of the perturbed positive adjoint modulo a
base pullback.  In the divisorial case that lemma makes its pushdown
rational Cartier.  These are therefore steps for the original data.

Finally Lemma~\ref{lem:negativity} proves the discrepancy comparisons.
For gklt, positivity is preserved.  For gdlt, a new zero-discrepancy
center cannot have its general point in an exceptional locus, by
strictness.  Its corresponding old zero center is generically in the
unchanged good open set.  The transformed boundary is SNC and the
b-data descend there.  This gives a gdlt good open set on the output.
The same hypotheses hold for the next relative morphism to $V$, so
the construction continues whenever relative nefness fails.
\end{proof}

\subsection{Finiteness on a compact relative base}

The next proposition constructs some finite relative program.  It will
be used to replace a single step of an arbitrary sequence by a finite
program on a higher model.  It does not assert termination of every
relative program.  We first prove the finite-model statement needed
for this construction.

\begin{lemma}[Local finiteness of marked ample models]\label{lem:marked-models}
Let $T\to V$ be projective and bimeromorphic and let
$D_0,\ldots,D_s$ be rational gklt adjoints on $T$ with a common
projective carrier and nef data over $V$.  On a smaller neighborhood
of a prescribed Stein compact, the rational adjoints in the simplex
of their convex combinations have only finitely many marked relative
ample models.  Here a marking is the bimeromorphic map from $T$,
and an ample model is required to have an effective exceptional
pullback difference.  No local $\Q$-factoriality of $T$ is assumed.
\end{lemma}

\begin{proof}
Lemma~\ref{lem:ordinary-representative} supplies ordinary klt
representatives with a common positive relatively ample part.
Choose a common projective log resolution $a\colon S\to T$ by
successive blowups and principalizations with centers exceptional over
$T$, resolving all fixed supports.  This choice can be made with an
effective exceptional divisor $F$ such that $-F$ is $a$-ample:
for each blowup take its effective tautological exceptional divisor,
and inductively add it to a sufficiently large positive multiple of the
pullback of the divisor for the preceding composite.  Relative
ampleness under composition proves the assertion; further such blowups
preserve it.  For the identity resolution take $F=0$.
This is the blowup-resolution construction used in
\cite[Theorem~2.13, Remark~2.14, and the proof of Corollary~3.7]{DHP2024}.
For each vertex choose every exceptional coefficient strictly between
the maximum of zero and its crepant coefficient and one.  We obtain
\[
 K_S+\Gamma_j=a^*D_j+E_j,
 \qquad E_j\geq0\text{ exceptional},\quad
 0\leq\Gamma_j<1.
\]
All supports, including that of $F$, are SNC.  The ample
part downstairs permits a common relatively ample part upstairs.
More explicitly, scale the chosen $F$ by a small positive rational number so that
$a^*H-F$ is ample over $V$, where $H$ is the fixed ample part
downstairs.  For a sufficiently small $\tau>0$ put
$A=\tau(a^*H-F)$.  The finitely many exceptional coefficients of the
$\Gamma_j$ have positive margins from both zero and one.  Then
$\Gamma_j-A$ is effective with coefficients below one; their
supports, including that of $F$, are SNC.  One can either use this
relatively ample rational divisor $A$ directly or replace it by a
divided general member and make the corresponding simultaneous
linear-equivalence changes in the boundaries.

The multigraded ring of the $K_S+\Gamma_j$ is locally finitely
generated by \cite[Theorem~3.1]{DHP2024}.  For nonnegative integers
$n_j$, the correction $\sum n_jE_j$ is effective and exceptional
over $T$.  Normality gives
$a_*\cO_S(\sum n_jE_j)=\cO_T$.  The projection formula therefore
identifies every multigraded piece with the corresponding piece for
the $D_j$, compatibly with multiplication.  Passing to one common
Veronese removes all rounding.  Thus their actual multigraded ring
is locally finitely generated as well.

For completeness, finite generation gives the required finiteness as
follows.  On a sufficiently small neighborhood choose homogeneous
generators with multidegrees $d_1,\ldots,d_N\in\N^{s+1}$.  For a
rational weight $w$ in the simplex, a point of the relative spectrum
is semistable for the diagonal ring of weight $w$ precisely when
$w$ is in the cone generated by those $d_i$ whose generators are
nonzero at the point.  To verify this equivalence, clear the rational
coefficients of such a cone expression and form the corresponding
monomial; conversely a nonvanishing homogeneous section contains a
nonvanishing monomial of that weight.  There are only finitely many
subsets of the generators and hence finitely many resulting open
sets.  The diagonal Proj is the quotient of its open set, obtained by
the degree-zero localizations of these monomials.  Uniqueness of this
quotient shows that weights with the same open set give the same
ample model.  All identifications commute with the map from $T$,
which is determined on the bimeromorphic isomorphism locus.  A finite
cover of the Stein compact makes the number of possibilities finite
there.  This is also the marked finite-model conclusion of
\cite[Theorem~E]{Fujino2022}, now with the ordinary representatives
and the smooth reduction justified.
\end{proof}

\begin{proposition}[A finite relative gklt run]\label{prop:relative-finite}
In the setting of Proposition~\ref{prop:relative-steps}, assume the
pair is gklt.  There is a finite sequence of elementary steps over
$V$ whose final adjoint is curve-nef over $V$.
\end{proposition}

\begin{proof}
Let $D=K_T+B+M_T$.  We may suppose $D$ is not already relatively
nef.

\smallskip
\noindent\emph{A scaling class valid on every later model.}
Choose global rational line bundles $L_1,\ldots,L_r$ whose
degrees form a basis of $N^1_{\mathrm{gl}}(T/V)_{\Q}$, and write
the class of $D$ as $d\in\Q^r$.

These bundles continue to span the degree space after every finite
sequence of elementary steps.  Indeed every line bundle on a later
model can be transported to a rational line bundle on the preceding
one by Lemma~\ref{lem:trace} and global strong $\Q$-factoriality.
For a divisorial contraction its pullback is also such a bundle;
exceptional-divisor adjustments are available on the source.
Moreover a bundle numerically trivial over $V$ remains so after a
step.  It first descends across that contraction by
Lemma~\ref{lem:descent}.  Any curve on the new side over $V$ has a
projective lift to a common model and then a curve upstairs mapping
onto it with positive degree.  Intersecting the pullbacks of the
descended bundle proves its zero degree.  Thus every later relative
curve determines a rational functional $v\in\Q^r$, nonzero because
the later space has a relatively ample line bundle.

Choose a real vector $h\in\R^r$ defining a real linear combination
$H$ of the chosen rational line bundles, with $H$ and $D+H$
relatively ample.  We choose it outside
the following countable family of proper rational hyperplanes:
\begin{align}
 v(h)&=0 &&(0\ne v\in\Q^r),\label{eq:scaling-generic1}\\
 v(d)w(h)-w(d)v(h)&=0
 &&\left(\begin{array}{l}v,w\in\Q^r\text{ independent},\\
                  (v(d),w(d))\ne(0,0)\end{array}\right).
 \label{eq:scaling-generic2}
\end{align}
The ample conditions define a nonempty open set, and a countable
union of proper hyperplanes cannot exhaust it.  The rational
functionals here range over all possibilities, so this single choice
works at every later model; no enumeration of the models is required.

\smallskip
\noindent\emph{The scaling steps and their strictly decreasing thresholds.}
On a current model, let $D_i$ be the reflexive transform of $D$,
and let $H_i$ be the same real linear combination of the reflexive
transforms of the rational line bundles defining $H$; across a
divisorial step these transforms are the corresponding pushforwards.
Regard $H$
also as nef b-data over $V$, carried on the initial space.  Its
pullback to a common carrier is relatively nef, so for every
$t\geq0$ the initial generalized pair with adjoint $D+tH$ is gklt.
At each preceding step of threshold $t_j$, its adjoint has negative
degree for $0\leq t<t_j$ and zero degree for $t=t_j$.
Lemma~\ref{lem:negativity} consequently shows that the transformed
pair is gklt throughout the parameter interval below the preceding
threshold.

Start with the relatively ample class $D+H$.  On a current model
take the lower endpoint $t_i$ of the nef segment extending downward
from the preceding parameter.  If $t_i=0$, then $D_i$ is nef and
the process stops.  Otherwise $J_i=D_i+t_iH_i$ is nef and not ample.
The local cone argument preceding Lemma~\ref{lem:sheaf-step}, applied
to this real gklt adjoint, shows that its zero face is generated by
finitely many curves.  A generating curve cannot have $D_i$-degree
zero, because then its $H_i$-degree would be zero, contrary to
\eqref{eq:scaling-generic1}.  Two independent such rays would violate
\eqref{eq:scaling-generic2}.  Hence the face is one ray in the global
degree space.  Its $D_i$-degree is negative: if it were positive,
lowering $t$ would be positive on this face, and the uniform cone
margin would give nefness below $t_i$, a contradiction.

Apply Proposition~\ref{prop:relative-steps} to this $D_i$-negative
ray.  The threshold class descends.  Here descent for a real class
requires no real-bundle base point free theorem: the perpendicular
hyperplane of the rational ray is rational, so express the threshold
class as a real combination of rational line-bundle classes of degree
zero and descend each by Lemma~\ref{lem:descent}.  The threshold
comparison is crepant and its transform is nef.

The thresholds strictly decrease.  In a small step, the zero face on
the positive side is again generated by curves by the real gklt
local cone argument.  Equations \eqref{eq:scaling-generic1}--
\eqref{eq:scaling-generic2} allow at most one ray.  Flipped curves
supply that ray and have positive $D$-degree.  Lowering the parameter
is therefore positive on it, and the uniform margin makes the class
ample for a nonempty interval below the old threshold.  In a
divisorial step there are no zero curves at the old threshold:
lifting one to the old space would make it lie on the contracted
zero ray while still mapping onto a curve.  The same margin again
gives an ample interval.  At every stage, the open interval between
successive thresholds consists of relatively ample classes: an
interior zero curve would have zero degree for two parameters, hence
zero $H_i$-degree, contrary to \eqref{eq:scaling-generic1}.

\smallskip
\noindent\emph{Finitely many marked ample models.}
It remains to exclude infinitely many such intervals.  Choose finitely
many rational relatively ample bundles $H^{(1)},\ldots,H^{(s)}$
near $H$, with $H$ in the interior of their convex hull and with
$D+H^{(j)}$ relatively ample.  Consider the fixed rational simplex
of weights with vertices
\begin{equation}\label{eq:scaling-polytope}
 D,\ D+H^{(1)},\ldots,D+H^{(s)}.
\end{equation}
At an interior parameter of any one of the alleged infinitely many
ample intervals, all preceding intersections for that parameter are
strictly negative and final ampleness is open.  Choose a nearby
rational point of \eqref{eq:scaling-polytope} retaining these finitely
many strict conditions.  The current model is its marked ample
model: on a common resolution its pullback difference from the
initial adjoint is effective and exceptional over the current model
by composition of the preceding comparisons.  The corresponding
relative section algebras agree, since pushing an effective
exceptional correction to a normal target gives the structure sheaf.

Each vertex in \eqref{eq:scaling-polytope} has nef b-data on the
initial model and is gklt.  Apply Lemma~\ref{lem:marked-models} on
each member of a finite Stein compact cover of $V$.  It gives finitely
many local marked ample models, hence finitely many possible tuples
of such local models.  Two global models giving the same tuple are
uniquely isomorphic over the interiors of these compacts: the
identification is fixed on the common bimeromorphic open set and
extends by the marked local isomorphism.  The identifications agree
on overlaps by normality, so they glue.  Thus only finitely many
global marked ample models occur.

A marked model cannot recur across a nontrivial intervening step.
For a fixed marking $T\dashrightarrow Y$, choose a common resolution.
The pullback difference for $D+tH$ is affine in $t$.  Its coefficient
inequalities of effectivity are convex, and the relatively ample cone
of $Y$ is convex.  The parameters for which $Y$ is its marked ample
model therefore form an interval.  If it occurred in two separated
ample intervals of the run, it would also be the ample model at every
intermediate ample parameter.  Uniqueness of the relative Proj would
identify each intermediate model with the same marking.  But adjacent
models across a nontrivial elementary step have different markings,
as their adjoint has opposite relative signs.  This is impossible.
Finitely many marked models therefore imply finitely many steps.
\end{proof}

\begin{remark}
The proof concerns some finite relative run over a bimeromorphic
compact base.  The finite-cover argument for ordinary pairs is also
explicit in \cite[Theorems~2.1 and~2.4]{FujinoCompact2026}.  We have
given the additional generalized-data, global-sheaf, and marked-model
arguments, because an ordinary theorem assuming $\Q$-factoriality near
a Stein compact does not by itself prove the statement above.
\end{remark}

\subsection{Extraction of prescribed low places}

\begin{proposition}[Exact extraction]\label{prop:extraction}
Let $(T,B+\bM)$ be a rational gklt pair on a globally strongly
$\Q$-factorial compact K\"ahler space, with projective nef carrier.
For any specified finite set $\mathcal E$ of exceptional prime
divisors over $T$ with
$a(E;T,B+\bM)\leq1$, there is a projective modification
$h\colon U\to T$ such that:
\begin{enumerate}[label=\textup{(\roman*)}]
\item $U$ is compact K\"ahler and globally strongly $\Q$-factorial;
\item the $h$-exceptional prime divisors are exactly $\mathcal E$;
\item the crepant generalized pair on $U$ is gklt with effective
boundary.
\end{enumerate}
\end{proposition}

\begin{proof}
The projective realization statement in Lemma~\ref{lem:low-places}
places the finitely many designated divisors on a projective log
resolution carrying the nef data.  Its proof uses only log-smooth
discrepancy calculations and stratum blowups, independently of
extraction.  Write this resolution as $p\colon W\to T$.
Keep every designated coefficient at its crepant value $1-a(E)$,
which lies in $[0,1)$.  Keep the nonexceptional coefficients as well.
For every other exceptional prime choose a coefficient in $[0,1)$
strictly greater than its crepant coefficient.  Since the pair is
gklt, there are only finitely many such choices and they are possible.
With this effective SNC boundary $\Gamma$ we have
\[
 K_W+\Gamma+M_W=p^*D_T+G,
 \qquad G\geq0,
\]
where the support of $G$ is exactly the set of undesired exceptional
primes.  The new pair is gklt.

Apply Proposition~\ref{prop:relative-finite} over $T$.  At every
stage the adjoint is the pullback of $D_T$ plus the effective
transform of $G$, and the latter remains exceptional over $T$.
A divisorial negative step cannot contract a prime outside its
support: a general covering contracted curve in that prime avoids
containment in the support and has nonnegative intersection with
the effective rational Cartier divisor $G$, whereas its adjoint
degree is negative.  Small steps contract no prime divisor.  Hence
every designated prime survives, and no nonexceptional prime of the
resolution is lost.

At the relative nef endpoint $h\colon U\to T$, the remaining
effective exceptional correction $G_U$ is $h$-nef.  Negativity gives
$G_U\leq0$, hence $G_U=0$.  Every undesired prime has therefore been
contracted.  The equality is now crepant, and the retained boundary
is effective.  The global strong condition, K\"ahlerness, and gklt
property follow at each step from Proposition~\ref{prop:relative-steps}.
The exceptional primes are exactly the specified ones, as claimed.
\end{proof}
